\documentclass[12pt]{article}
\usepackage[utf8]{inputenc}
\usepackage[T1]{fontenc}
\usepackage[margin=1in]{geometry}
\usepackage{setspace}
\usepackage{xcolor}
\usepackage[colorlinks=true, urlcolor=blue!60!black, linkcolor=blue!60!black]{hyperref}
\setlength{\parindent}{2em}
\setlength{\parskip}{1.25em}
\renewcommand{\baselinestretch}{1.0}
\title{\textit{Shelby County v. Holder}, 570 U.S. 529 (2013)}
\author{Case Brief}
\date{}
\begin{document}
\maketitle

% Case Brief: Shelby County v. Holder, 570 U.S. 529 (2013)
% Constitutional Law Briefing Book — Chapter: Reconstruction Amendments & State Action

\section{Case Metadata}

\begin{tabular}{@{}p{0.26\textwidth}p{0.70\textwidth}@{}}
\textbf{Case Name:} & \emph{Shelby County, Alabama v. Eric H. Holder, Jr., Attorney General, et al.} \\
\textbf{Citation:} & 570 U.S. 529 (2013); 133 S. Ct. 2612; 186 L. Ed. 2d 651 (No. 12--96) \\
\textbf{Court \& Dates:} & Supreme Court of the United States; Argued February 27, 2013; Decided June 25, 2013 \\
\textbf{Lower Courts:} & \emph{Shelby County v. Holder}, 811 F. Supp. 2d 424 (D.D.C. 2011) (upholding \S\S~4(b) and 5), \emph{aff'd}, 679 F.3d 848 (D.C. Cir. 2012) (2--1, with Judge Stephen Williams dissenting), \emph{cert. granted}, 568 U.S. 1006 (2012) \\
\textbf{Vote:} & 5--4, reversing the D.C. Circuit \\
\textbf{Majority Opinion:} & Chief Justice Roberts, joined by Justices Scalia, Kennedy, Thomas, and Alito \\
\textbf{Concurring Opinion:} & Justice Thomas (joining the Court's opinion in full but writing separately to argue that \S~5 is also unconstitutional) \\
\textbf{Dissenting Opinion:} & Justice Ginsburg, joined by Justices Breyer, Sotomayor, and Kagan \\
\textbf{Constitutional \& Statutory Provisions:} & U.S. Const. amend. X; U.S. Const. amend. XIV, \S\S~1, 5; U.S. Const. amend. XV, \S\S~1, 2; Voting Rights Act of 1965 (VRA), \S~2 (formerly 42 U.S.C. \S~1973, now codified at 52 U.S.C. \S~10301), \S~4(b) (formerly 42 U.S.C. \S~1973b(b), now 52 U.S.C. \S~10303(b)), and \S~5 (formerly 42 U.S.C. \S~1973c, now 52 U.S.C. \S~10304) \\
\textbf{Primary Sources:} & \href{https://www.law.cornell.edu/supremecourt/text/12-96}{Cornell LII Full Opinion}; \href{https://supreme.justia.com/cases/federal/us/570/529/}{Justia U.S. Supreme Court Center}; \href{https://tile.loc.gov/storage-services/service/ll/usrep/usrep570/usrep570529/usrep570529.pdf}{U.S. Reports (Library of Congress PDF)} \\
\end{tabular}

\section{Facts and Procedural History}

\subsection{Statutory Architecture of the Voting Rights Act of 1965}
Ratified in 1870, the Fifteenth Amendment commands that ``[t]he right of citizens of the United States to vote shall not be denied or abridged by the United States or by any State on account of race, color, or previous condition of servitude,'' and empowers Congress ``to enforce this article by appropriate legislation.'' For nearly a century, however, southern States systematically nullified that guarantee through literacy tests, good-moral-character requirements, voucher rules, poll taxes, white primaries, and defiant evasion of federal court decrees (\href{https://www.law.cornell.edu/supremecourt/text/12-96}{Cornell LII, Roberts Op. Part I-A}).

Congress responded in 1965 with the Voting Rights Act (VRA), pairing permanent nationwide protections with temporary, geographically targeted ``extraordinary measures'':
\begin{itemize}
    \item \textbf{Section 2 (Permanent, Nationwide):} Originally prohibited any voting qualification, prerequisite, standard, practice, or procedure imposed to deny or abridge the right to vote on account of race or color. As amended in 1982, \S~2 prohibits any voting practice that ``results in'' such a denial or abridgement nationwide, enforceable by both the federal government and private litigants. Section 2 was not challenged in \emph{Shelby County}.
    \item \textbf{Section 4(b) (The Coverage Formula):} Identified ``covered jurisdictions''---States or political subdivisions that, as of November 1, 1964, maintained a ``test or device'' as a prerequisite to voting and had less than 50\% voter registration or turnout in the 1964 presidential election. Initially, \S~4(b) covered six whole States (Alabama, Georgia, Louisiana, Mississippi, South Carolina, and Virginia) plus 39 counties in North Carolina and one in Arizona. Under \S~4(a), a covered jurisdiction could ``bail out'' by proving a clean record over a statutory lookback period.
    \item \textbf{Section 5 (Federal Preclearance):} Barred covered jurisdictions from implementing \emph{any} change in voting laws or procedures---from statewide redistricting down to moving a local polling place---until approved (``precleared'') either by the U.S. Attorney General (within 60 days, extendable upon requests for more information) or by a three-judge U.S. District Court in the District of Columbia. The jurisdiction bore the burden of proving that the proposed change had neither the purpose nor the effect of denying or abridging the right to vote on account of race or color.
\end{itemize}

\subsection{Reauthorizations and Prior Supreme Court Precedent}
Sections 4 and 5 were originally enacted to expire after five years. In \emph{South Carolina v. Katzenbach}, 383 U.S. 301 (1966), the Supreme Court upheld both provisions as ``rational in both practice and theory'' given ``exceptional conditions'' of pervasive voting discrimination. Congress subsequently reauthorized the temporary provisions four times:
\begin{itemize}
    \item \textbf{1970 (5-year extension):} Updated \S~4(b) to incorporate 1968 election data (bringing in counties in California, New Hampshire, and New York) and extended the ban on tests and devices nationwide. Upheld in \emph{Georgia v. United States}, 411 U.S. 526 (1973).
    \item \textbf{1975 (7-year extension):} Updated \S~4(b) to incorporate 1972 election data, expanded ``test or device'' to include English-only voting materials in jurisdictions with language-minority populations exceeding 5\% (bringing in Alaska, Arizona, Texas, and additional counties in California, Florida, Michigan, New York, North Carolina, and South Dakota), and made the nationwide ban on tests and devices permanent. Upheld in \emph{City of Rome v. United States}, 446 U.S. 156 (1980).
    \item \textbf{1982 (25-year extension):} Retained the 1975 coverage formula unchanged while liberalizing bailout to allow political subdivisions within covered States to seek bailout after ten years of compliance. Upheld in \emph{Lopez v. Monterey County}, 525 U.S. 266 (1999).
    \item \textbf{2006 (25-year extension, through 2031):} Following 21 congressional hearings and a legislative record exceeding 15,000 pages, Congress passed the Fannie Lou Hamer, Rosa Parks, and Coretta Scott King Voting Rights Act Reauthorization and Amendments Act of 2006 by votes of 390--33 in the House and 98--0 in the Senate, signed by President George W. Bush. Congress \emph{left \S~4(b)'s 1964/1968/1972 coverage formula untouched} while expanding \S~5's substantive standard to override \emph{Reno v. Bossier Parish School Board}, 528 U.S. 320 (2000) (\emph{Bossier II}), and \emph{Georgia v. Ashcroft}, 539 U.S. 461 (2003), prohibiting voting changes with ``any discriminatory purpose'' or that diminish minority voters' ability ``to elect their preferred candidates of choice'' (\href{https://www.law.cornell.edu/supremecourt/text/12-96}{Cornell LII, Roberts Op. Part I-A; Ginsburg Dissent Part I}).
\end{itemize}

In \emph{Northwest Austin Municipal Utility District No. One v. Holder}, 557 U.S. 193 (2009), an 8--1 Court (per Chief Justice Roberts, with Justice Thomas concurring in the judgment in part and dissenting in part) avoided deciding the constitutionality of \S~5 by construing the statutory bailout provision broadly to allow a Texas utility district to seek bailout. Crucially, however, \emph{Northwest Austin} warned that ``[t]hings have changed in the South,'' that the VRA's ``current burdens . . . must be justified by current needs,'' and that a departure from ``the fundamental principle of equal sovereignty'' requires a showing that disparate geographic coverage is ``sufficiently related to the problem that it targets'' (\href{https://www.law.cornell.edu/supremecourt/text/12-96}{Cornell LII, Roberts Op. Part I-A}).

\subsection{Litigation Below}
Shelby County, Alabama---located in a covered State and ineligible for bailout due to recent DOJ preclearance objections within the county (including the 2008 Calera redistricting dispute)---sued Attorney General Eric Holder in the U.S. District Court for the District of Columbia in 2010, seeking a declaratory judgment that \S\S~4(b) and 5 were facially unconstitutional and a permanent injunction against their enforcement.

The District Court granted summary judgment to the Attorney General, holding that the 2006 congressional record adequately justified reauthorizing both \S~5 and \S~4(b). \emph{Shelby County v. Holder}, 811 F. Supp. 2d 424, 508 (D.D.C. 2011). A divided panel of the D.C. Circuit affirmed, concluding that six categories of record evidence (DOJ objections, requests for more information, successful \S~2 suits, federal observers, \S~5 enforcement suits, and deterrence) supported \S~5, and that---although \S~4(b) presented ``a close question''---disparate rates of successful \S~2 litigation and \S~5's deterrent effect justified retaining the coverage formula. 679 F.3d 848, 862--63, 879--83 (D.C. Cir. 2012). Judge Stephen Williams dissented, finding no positive correlation between \S~4(b) coverage and low Black registration or turnout, noting that the five worst uncovered States had worse \S~2 records than eight covered States, and concluding that \S~4(b) was ``irrational'' (\href{https://www.law.cornell.edu/supremecourt/text/12-96}{Cornell LII, Roberts Op. Part I-B}). The Supreme Court granted certiorari.

\section{Issue(s) Presented}

Whether Congress's 2006 reauthorization of the Voting Rights Act's \S~4(b) coverage formula---which continued to subject nine States and several additional local subdivisions to \S~5 federal preclearance based on voting tests and voter registration or turnout data from 1964, 1968, and 1972---exceeded Congress's enforcement authority under the Fourteenth and Fifteenth Amendments in light of the Tenth Amendment and the constitutional principle of the equal sovereignty of the States.

\section{Holding}

Yes. In a 5--4 decision authored by Chief Justice Roberts, the Supreme Court held that \textbf{\S~4(b) of the Voting Rights Act is unconstitutional} because its coverage formula is based on 40-year-old data and eradicated practices that no longer speak to current conditions, and therefore can no longer be used as a basis for subjecting jurisdictions to federal preclearance under \S~5 (\href{https://www.law.cornell.edu/supremecourt/text/12-96}{Cornell LII, Roberts Op. Part III-D}).

The Court expressly declined to hold \S~5 itself unconstitutional and emphasized that its ruling ``in no way affects the permanent, nationwide ban on racial discrimination in voting found in \S~2.'' Practically, however, invalidating \S~4(b) rendered \S~5's preclearance regime inoperable nationwide unless and until Congress enacts a new coverage formula grounded in current conditions (aside from targeted judicial ``bail-in'' orders under \S~3(c), 52 U.S.C. \S~10302(c)).

\section{Reasoning}

Chief Justice Roberts's majority opinion rests on three interlocking doctrinal pillars: structural federalism and the ``equal sovereignty'' of the States, empirical transformation since 1965, and the irrationality of re-enacting a 1965-era coverage formula in 2006.

\subsection{Federalism, State Election Autonomy, and Equal Sovereignty}
Although the Supremacy Clause prevents state laws from contravening valid federal law, the Framers expressly rejected James Madison's proposal at the Constitutional Convention to give the Federal Government a general ``negative'' (veto) over state enactments before they take effect. Outside the Supremacy Clause, the Tenth Amendment reserves to the States broad autonomy, including ``the power to regulate elections'' (\emph{Gregory v. Ashcroft}, 501 U.S. 452, 461--62 (1991)). Moreover, the Court elevated what \emph{Northwest Austin} called the ``fundamental principle of equal sovereignty'' among the States---drawing on \emph{Coyle v. Smith}, 221 U.S. 559, 567, 580 (1911) (``a union of States, equal in power, dignity and authority''). While acknowledging that \emph{Katzenbach} rejected equal sovereignty as an absolute bar outside the admission of new States, the majority held that the principle ``remains highly pertinent in assessing subsequent disparate treatment of States'' (\href{https://www.law.cornell.edu/supremecourt/text/12-96}{Cornell LII, Roberts Op. Part II-A}).

Section 5 and \S~4(b) sharply depart from both principles: covered States must ``beseech the Federal Government for permission'' to enact even innocuous voting changes, shifting the burden of proof onto the State and imposing delays of months or years, while neighboring uncovered States implement identical laws immediately. In 1966, \emph{Katzenbach} tolerated this ``strong medicine'' solely because ``exceptional conditions'' of pervasive, entrenched disenfranchisement made the 1965 formula ``rational in both practice and theory'' by tying cause (discriminatory tests) to effect (low registration and turnout) (\href{https://www.law.cornell.edu/supremecourt/text/12-96}{Cornell LII, Roberts Op. Part II-B}).

\subsection{Dramatic Empirical Progress Since 1965}
The majority emphasized that ``[n]early 50 years later, things have changed dramatically''---largely because the VRA worked. Literacy tests and similar devices had been banned nationwide for over 40 years. Relying on Congress's own 2006 committee reports and 2012 Census Bureau data, Chief Justice Roberts highlighted a sharp table comparing 1965 and 2004 Black-white voter registration gaps in the six originally covered States (e.g., Mississippi's gap narrowed from 63.2 percentage points to $-3.8$, with Black registration exceeding white registration; Alabama's gap narrowed from 49.9 points to 0.9 points). By 2012, Black voter turnout exceeded white turnout in five of the six originally covered States, with a gap under 0.5\% in the sixth; Black elected officials in those States had increased by roughly 1,000\%; and DOJ preclearance objection rates fell from 14.2\% in the first decade to 0.16\% in the decade preceding 2006 (\href{https://www.law.cornell.edu/supremecourt/text/12-96}{Cornell LII, Roberts Op. Part II-C}).

Yet rather than easing \S~5 or narrowing \S~4(b), Congress in 2006 extended preclearance for another 25 years and tightened \S~5's substantive bar. The majority also rejected the government's deterrence argument as non-falsifiable: ``no matter how `clean' the record of covered jurisdictions, the argument could always be made that it was deterrence that accounted for the good behavior,'' effectively immunizing \S~5 from judicial review (\href{https://www.law.cornell.edu/supremecourt/text/12-96}{Cornell LII, Roberts Op. Part II-C}).

\subsection{Irrationality of the 1965 Formula Under Current Conditions}
Applying \emph{Northwest Austin}'s dual mandate---that ``current burdens'' must be justified by ``current needs,'' and that disparate geographic coverage must be ``sufficiently related to the problem that it targets''---the majority invalidated \S~4(b) on three grounds:
\begin{enumerate}
    \item \textbf{Rejection of ``Reverse-Engineering'':} The government argued that Congress identified the jurisdictions it wanted covered and used \S~4(b) simply as a proxy to describe them. The majority held that \emph{Katzenbach} required the formula's criteria themselves to be relevant to the targeted evil; failing to establish even relevance is fatal when singling out a disfavored subset of States (\href{https://www.law.cornell.edu/supremecourt/text/12-96}{Cornell LII, Roberts Op. Part III-B}).
    \item \textbf{``History Did Not End in 1965'':} The Fifteenth Amendment ``is not designed to punish for the past; its purpose is to ensure a better future.'' Congress could not justify treating States differently in 2006 solely because they had different histories in 1965. If Congress had drafted \S~4(b) from scratch in 2006 using 40-year-old data and long-banned tests, the formula would have been plainly irrational (\href{https://www.law.cornell.edu/supremecourt/text/12-96}{Cornell LII, Roberts Op. Parts III-B, III-D}).
    \item \textbf{Disconnect Between the 2006 Record and the Statutory Formula:} Even accepting that Congress compiled over 15,000 pages of evidence regarding ``second-generation barriers'' (such as vote dilution and racial gerrymandering), Congress \emph{did not use that record to craft a new coverage formula}. Worse, pointing to second-generation vote-dilution barriers only underscored the irrationality of a formula still tethered to first-generation ballot-access tests. And Shelby County could bring a facial challenge because it was subjected to preclearance via an unconstitutional selection rule---analogous to ``a driver pulled over pursuant to a policy of stopping all redheads'' who may challenge the policy even if his license happens to be expired (\href{https://www.law.cornell.edu/supremecourt/text/12-96}{Cornell LII, Roberts Op. Part III-C}).
\end{enumerate}

Notably, the majority cited \emph{Northwest Austin} to confirm that its analysis governed under both the Fourteenth and Fifteenth Amendments (n.1), without expressly deciding whether Fifteenth Amendment enforcement triggers \emph{Katzenbach}'s rational-means standard or \emph{City of Boerne v. Flores}'s ``congruence and proportionality'' standard.

\section{Concurring and Dissenting Opinions}

\subsection{Concurring Opinion of Justice Thomas}
Justice Thomas joined the Court's opinion in full but wrote separately to argue that \textbf{\S~5 itself is unconstitutional}. In his view, the majority's own reasoning---that the 2006 record did not show anything approaching the ``pervasive,'' ``flagrant,'' ``widespread,'' and ``rampant'' discrimination of 1965, even as Congress raised \S~5's substantive bar---compellingly proved that preclearance could no longer satisfy ``current burdens'' justified by ``current needs.'' Leaving \S~5 formally intact, Thomas argued, ``needlessly prolongs the demise of that provision'' (\href{https://www.law.cornell.edu/supremecourt/text/12-96}{Cornell LII, Thomas Concurrence}).

\subsection{Dissenting Opinion of Justice Ginsburg (Joined by Justices Breyer, Sotomayor, and Kagan)}
Justice Ginsburg delivered a forceful, institutionalist dissent framing the case as a question of \emph{who decides} whether preclearance remains necessary: the Supreme Court, or Congress---the branch expressly charged by the Reconstruction Amendments with enforcing voting equality ``by appropriate legislation'' (\href{https://www.law.cornell.edu/supremecourt/text/12-96}{Cornell LII, Ginsburg Dissent Intro}).

\begin{itemize}
    \item \textbf{Broad Congressional Enforcement Power:} Invoking \emph{McCulloch v. Maryland}, 4 Wheat. 316, 421 (1819), and \emph{Katzenbach}'s rational-means standard, Ginsburg argued that congressional power is at its zenith when targeting racial discrimination in the right ``preservative of all rights'' (\emph{Yick Wo v. Hopkins}). Unlike the First Amendment's restriction (``Congress shall make no law''), the Reconstruction Amendments were transformative affirmative grants of legislative authority. Moreover, reauthorization of a previously upheld, functioning statute deserves heightened deference; demanding a record as stark as 1965 creates an impossible catch-22 where a statute's success becomes the legal basis for destroying it (\href{https://www.law.cornell.edu/supremecourt/text/12-96}{Cornell LII, Ginsburg Dissent Part II}).
    \item \textbf{Second-Generation Barriers and the 2006 Record:} Voting discrimination did not disappear; like ``battling the Hydra,'' it evolved from first-generation ballot-access denials into ``second-generation barriers'' (at-large elections, racial gerrymandering, discriminatory annexations, polling-place relocations, and early-voting cuts). Between 1982 and 2006, DOJ blocked over 700 discriminatory changes---more objections (626 from 1982--2004) than in 1965--1982 (490)---and over 800 proposals were modified or withdrawn after DOJ inquiries. Furthermore, the Katz study showed that covered jurisdictions (under 25\% of U.S. population) accounted for 56\% of successful \S~2 lawsuits, nearly four times the rate of noncovered jurisdictions per capita, alongside sharper racial polarization. Combined with dynamic ``bail-out'' (nearly 200 jurisdictions successfully exited) and ``bail-in'' provisions, the formula remained well calibrated (\href{https://www.law.cornell.edu/supremecourt/text/12-96}{Cornell LII, Ginsburg Dissent Part III}).
    \item \textbf{Three ``Disturbing Lapses'' by the Majority:} First, the majority ignored \emph{United States v. Salerno}, 481 U.S. 739 (1987), and the VRA's broad severability clause (42 U.S.C. \S~1973p) by entertaining a facial challenge from Shelby County, Alabama---a jurisdiction whose own recent record (including the 2008 Calera defiance and 2010 FBI recordings of Alabama legislators in \emph{United States v. McGregor}) plainly justified preclearance as applied; ``[h]ubris is a fit word for today's demolition of the VRA.'' Second, the majority silently overruled \emph{Katzenbach}'s explicit holding that equal sovereignty ``applies only to the terms upon which States are admitted to the Union,'' elevating \emph{Northwest Austin} dictum without a word of \emph{stare decisis} analysis and threatening numerous federal statutes that treat States differently (e.g., PASPA sports-gambling grandfathering, Yucca Mountain nuclear waste rules). Third, the majority ignored the 15,000-page legislative record and dismissed the risk of backsliding---memorably declaring that ``[t]hrowing out preclearance when it has worked and is continuing to work to stop discriminatory changes is like throwing away your umbrella in a rainstorm because you are not getting wet'' (\href{https://www.law.cornell.edu/supremecourt/text/12-96}{Cornell LII, Ginsburg Dissent Part IV}).
\end{itemize}

\section{Key Quotations}

All quotations below are verified verbatim against the primary opinion text (\href{https://www.law.cornell.edu/supremecourt/text/12-96}{Cornell LII}; \href{https://tile.loc.gov/storage-services/service/ll/usrep/usrep570/usrep570529/usrep570529.pdf}{570 U.S. 529 (2013)}):

\begin{enumerate}
    \item \textbf{Extraordinary Measures and Equal Sovereignty (Chief Justice Roberts, 570 U.S. at 534--35):}
    \begin{quote}
    ``The Voting Rights Act of 1965 employed extraordinary measures to address an extraordinary problem. Section 5 of the Act required States to obtain federal permission before enacting any law related to voting---a drastic departure from basic principles of federalism. And \S~4 of the Act applied that requirement only to some States---an equally dramatic departure from the principle that all States enjoy equal sovereignty.''
    \end{quote}

    \item \textbf{The Governing Doctrinal Standard (Chief Justice Roberts, 570 U.S. at 542):}
    \begin{quote}
    ``In \emph{Northwest Austin}, we stated that `the Act imposes current burdens and must be justified by current needs.' 557 U. S., at 203. And we concluded that `a departure from the fundamental principle of equal sovereignty requires a showing that a statute's disparate geographic coverage is sufficiently related to the problem that it targets.' \emph{Ibid.}''
    \end{quote}

    \item \textbf{``History Did Not End in 1965'' (Chief Justice Roberts, 570 U.S. at 552--53):}
    \begin{quote}
    ``But history did not end in 1965. By the time the Act was reauthorized in 2006, there had been 40 more years of it. \ldots{} The Fifteenth Amendment commands that the right to vote shall not be denied or abridged on account of race or color, and it gives Congress the power to enforce that command. The Amendment is not designed to punish for the past; its purpose is to ensure a better future.''
    \end{quote}

    \item \textbf{Scope of Holding and Invitation to Congress (Chief Justice Roberts, 570 U.S. at 557):}
    \begin{quote}
    ``Our decision in no way affects the permanent, nationwide ban on racial discrimination in voting found in \S~2. We issue no holding on \S~5 itself, only on the coverage formula. Congress may draft another formula based on current conditions. \ldots{} Our country has changed, and while any racial discrimination in voting is too much, Congress must ensure that the legislation it passes to remedy that problem speaks to current conditions.''
    \end{quote}

    \item \textbf{``Umbrella in a Rainstorm'' (Justice Ginsburg, Dissenting, 570 U.S. at 590):}
    \begin{quote}
    ``Throwing out preclearance when it has worked and is continuing to work to stop discriminatory changes is like throwing away your umbrella in a rainstorm because you are not getting wet.''
    \end{quote}
\end{enumerate}

\section{Doctrinal Significance and Modern Legacy (Through 2026)}

Although \emph{Shelby County} formally left \S~5 intact, invalidating \S~4(b) immediately terminated federal preclearance across all nine covered States and dozens of covered local subdivisions. Formerly covered States swiftly enacted voter-ID laws, early-voting cuts, and polling-place closures previously deterred by \S~5, while Congress has never passed a replacement formula (repeated versions of the John R. Lewis Voting Rights Advancement Act, including H.R.~4 in 2021 and S.~2523 in 2025--2026, stalled in the Senate). Several States (e.g., New York, California, Virginia) responded with state-level Voting Rights Acts.

With \S~5 immobilized, enforcement shifted to \S~2---which the Court progressively narrowed through 2026:
\begin{itemize}
    \item \textbf{\emph{Brnovich v. Democratic National Committee}, 594 U.S. 647 (2021):} A 6--3 Court (per Justice Alito) announced five ``guideposts'' emphasizing ``usual burdens of voting'' and anti-fraud interests, sharply limiting \S~2 vote-denial claims (\href{https://supreme.justia.com/cases/federal/us/594/19-1257/}{Justia: \emph{Brnovich}}).
    \item \textbf{\emph{Allen v. Milligan}, 599 U.S. 1 (2023):} A 5--4 Court (per Chief Justice Roberts) temporarily preserved \emph{Thornburg v. Gingles}, 478 U.S. 30 (1986), rejecting Alabama's ``race-neutral benchmark,'' but Justice Kavanaugh's concurrence echoed \emph{Shelby County} by warning that race-based redistricting under \S~2 ``cannot extend indefinitely into the future'' (\href{https://www.law.cornell.edu/supremecourt/text/21-1086}{Cornell LII: \emph{Milligan} (2023)}).
    \item \textbf{\emph{Louisiana v. Callais}, 608 U.S. 85 (2026) (decided Apr. 29, 2026):} Following October 2025 reargument, a 6--3 Court (per Justice Alito) struck down Louisiana's remedial congressional map (SB8) as an unconstitutional racial gerrymander and held that---to remain congruent and proportional under the Fifteenth Amendment---``\S~2 imposes liability only when the evidence supports a strong inference'' of intentional racial discrimination (\href{https://supreme.justia.com/cases/federal/us/608/24-109/}{Justia: \emph{Callais}}; \href{https://www.scotusblog.com/2026/04/in-major-voting-rights-act-case-supreme-court-strikes-down-redistricting-map-challenged-as-racia/}{SCOTUSblog}). Quoting \emph{Shelby County}'s ``things have changed dramatically'' rationale, \emph{Callais} updated \emph{Gingles}: illustrative maps cannot use race and must satisfy all state political/incumbency goals; plaintiffs must control for party affiliation under \emph{Rucho v. Common Cause}; and the totality inquiry focuses on present-day intentional discrimination. Justice Kagan's dissent called the ruling the ``now-completed demolition'' of the VRA. Five weeks later, on June 2, 2026, the Court relied on \emph{Callais} to grant Alabama interim relief in \emph{Allen v. Milligan}, 608 U.S. \_\_\_ (2026) (\href{https://www.law.cornell.edu/supremecourt/text/25A1314}{Cornell LII: \emph{Milligan} (2026)}).
\end{itemize}

\section{Pedagogical Critiques and Exam Hypotheticals}

\subsection{Key Pedagogical Critiques}
Critics note that ``equal sovereignty'' appears nowhere in the Constitution's text, that the Reconstruction Amendments expressly altered antebellum federalism, and that invalidating \S~4(b) facially ignored both \emph{Salerno} and the VRA's severability clause (\S~1973p) despite Alabama's own recent record (\href{https://www.law.cornell.edu/supremecourt/text/12-96}{Cornell LII: \emph{Shelby County}}). Moreover, \emph{Shelby County} left unresolved whether Fifteenth Amendment laws face \emph{Katzenbach}'s rational-means test or \emph{City of Boerne}'s congruence-and-proportionality test---an ambiguity \emph{Callais} later resolved in favor of \emph{Boerne}.

\subsection{Exam Hypothetical}
Suppose Congress enacts a nationwide, rolling 15-year preclearance formula triggered whenever any State incurs five or more final federal judgments of intentional Fourteenth or Fifteenth Amendment voting violations, with expedited 5-year bailout. \emph{Analysis:} Unlike the 1965 formula in \emph{Shelby County}, this trigger uses current data, tracks adjudicated constitutional violations, and treats all States symmetrically on its face, strongly satisfying equal sovereignty. A challenger's best counterargument---echoing Justice Thomas's concurrence---is that prior federal restraint under \S~5 remains unconstitutional absent \emph{Katzenbach}-level systemic defiance.

\subsection{``Opposite Side'' Contrast}
\begin{itemize}
    \item \textbf{Federal Enforcement / Deference (\emph{Katzenbach}; Ginsburg Dissent; Kagan Dissent in \emph{Callais}):} The Reconstruction Amendments purposefully subordinated state sovereignty to protect racial voting equality; courts should defer to Congress's 15,000-page record and prophylactic remedies rather than discarding ``your umbrella in a rainstorm.''
    \item \textbf{Equal Sovereignty / Colorblindness (Roberts Majority \& Thomas Concurrence; Alito Majority in \emph{Callais}):} Emergency departures from state equality and race neutrality cannot endure after conditions transform; the Fifteenth Amendment enforces current constitutional rights rather than punishing the past or policing nonjusticiable partisan gerrymanders.
\end{itemize}

\end{document}
